\sscp{Distribution of the blowgun}{15-04-03-03}
Although not a linguistic feature,
the demarcation in the distribution of blowguns also occurs near
the middle of our target region,
as shown in \frf{15-07}.
They are limited to western parts of Wallacea,
but after a gap are also found in parts of Oceania.
Blowguns are found continuously to the west
and north of the Blust line,
they straddle the line,
but they are found only sporadically to the east away from the
line, where there are a few isolated occurrences.

\begin{figure}[h]\centering
	\includegraphics[width=\columnwidth]{blowgun.pdf}
	\fcp{Distribution of blowguns \citep{Jett1970,Jett1991,Blench2009}}{15-07}
\end{figure}

Buru \it{supit} ‘blowgun’ reflects a PMP form without a medial
*NC cluster (see \frf{15-06}).
\citet{BlustTrussel2020} have *sumpit ‘blowpipe, blowgun’.
\citet[158]{Dempwolff1938} reconstructed *su(m)pit ‘narrow,
blowpipe’ with an optional nasal.
The final \it{t{\#}} in Buru could be either a retention or a productive nominal suffix \it{-t}.
The ACD form with the medial *mp cluster should yield an unattested
*[subit] or *[sobit] in Buru,
so the form without a medial nasal is necessary.
This is an example of a cultural word that disappears for extra-linguistic
reasons, probably coming into contact with a more efficient bow and arrow.
Spears are the weapon of choice on Buru for hunting anything
larger than a bird and for fighting.
There are no bows
and arrows on Buru.
Bows and arrows are associated in Buru folklore with Seram across the Manipa Strait.
The Buru term \it{foros-banat} ‘bow
and arrow’ does not follow regular sound correspondences (cf.
PMP *busuR ‘hunting bow’,
*panaq ‘shoot with bow’)
and is surely a loan.

\begin{figure}[h]\centering
	\frame{\includegraphics[width=0.5\columnwidth]{15-06.jpg}}
	\fcp{Blowgun on Buru}{15-06}
\end{figure}

\citet[49]{Blench2009} says:

\begin{quote}
“One of the most striking examples
of a technology spread by Austronesian contact is the blow-gun.
Blow-guns used for hunting or warfare may seem like a technology
that could be invented many times,
but their worldwide distribution is very patchy.
Entirely absent from Africa
and Western Eurasia, they are typical of SE Asia
and a distinctive zone between North
and South America \citep{Jett1970,Jett1991}.
Virtually all occurrences within the Old
and New World distributions are contiguous,
suggesting the technology was only adopted once
and diffused, rather than developing through convergent evolution.
In the Old World,
the blow-gun maps extremely well against the proposed Austronesian
migration sites listed in this paper,
primarily in the Austronesian-speaking areas of insular SE Asia
and the Pacific.”
\end{quote}

There is a gap between eastern Indonesia
and the islands of New Britain
and New Ireland near eastern New Guinea -- the region often associated
with the origins of Oceanic languages.
This is one more feature found in our target region that does
not align with other features.
